\section{A bounded auxiliary oracle under conditional sub-Gaussian noise}
\label{sec:bounded-noise-oracle}
The next lemma supplies the exact leading-one comparison used in the
optimistic raw theorem. It concerns a proof-only bounded prior; the raw
learner itself is not changed.

\begin{lemma}[Conditional-noise leave-one-out bound]
\label{lem:bounded-noise-oracle}
Let a prior $\pi$ be supported on functions with values in $[-M,M]$.
Suppose $Y=f_0(X)+\xi$, $|f_0|\le M$, and conditionally on $X$ the
centered noise $\xi$ is $\sigma^2$-sub-Gaussian. At any fixed
$T\ge256(M^2+\sigma^2)$, the ordinary all-data posterior mean satisfies
\[
 \E R(\bar f_{S_n})\le\rho R+\frac{T\KL(\rho\|\pi)}{n+1}
\]
for every deterministic comparison law $\rho\ll\pi$.
\end{lemma}
\begin{proof}
Rescale by $M$; it suffices to prove the statement at $M=1$.
For a fixed law $q$ on $U\in[-1,1]$, write its mean and variance as
$\mu,v$, and let $p_y$ be its tilt by $e^{-(U-y)^2/T}$.
The covariance identity used in the bounded finite-variance lemma gives
\[
 \E_{p_y}(U-y)^2\ge(\mu-y)^2+
 \left[1-\frac{4(1+|y|)^2}{T}e^{4(1+|y|)/T}\right]v.
\]
If $Y=a+\xi$, $|a|\le1$, and $\xi$ is centered sub-Gaussian, then
$\E e^{t|\xi|}\le2e^{\sigma^2t^2/2}$ and
$\E\xi^4\le16\sigma^4$. Consequently Cauchy--Schwarz gives, at
$T\ge256(1+\sigma^2)$,
\begin{align*}
 \E[(1+|Y|)^2e^{4(1+|Y|)/T}]
 &\le e^{8/T}\left[16e^{8\sigma^2/T^2}
             +8\sqrt2\,\sigma^2e^{16\sigma^2/T^2}\right]\\
 &\le32(1+\sigma^2).
\end{align*}
The expected coefficient of $v$ is therefore nonnegative. Thus
$\E_Y(\mu-Y)^2\le\E_Y\E_{p_Y}(U-Y)^2$.

Introduce $n+1$ IID observations, using exactly the same prior and
 temperature as for the $n$-sample learner. Conditional on the other
$n$ observations and $X_i$, apply this inequality to the predictive law
of the posterior omitting observation $i$. Its update by observation
$i$ is the full $(n+1)$-sample posterior. Averaging over $i$ bounds the
expected test risk of the ordinary $n$-sample mean by the expected
training loss of that full posterior. Its Gibbs variational inequality
bounds the latter by $\rho R+T\KL(\rho\|\pi)/(n+1)$.
All integrals are finite because predictions are bounded and labels
have finite second moment. Rescaling proves the stated lemma.
\end{proof}
